\documentclass[11pt]{article}

\usepackage[margin=1in]{geometry}
\usepackage{amsmath,amssymb,amsfonts,amsthm,mathtools,bm}
\usepackage{algorithm}
\usepackage{algpseudocode}
\usepackage{booktabs}
\usepackage[numbers,sort&compress]{natbib}
\usepackage{microtype}
\usepackage{xcolor}
\usepackage[hidelinks]{hyperref}
\usepackage[nameinlink,capitalise,noabbrev]{cleveref}

\newtheorem{theorem}{Theorem}
\newtheorem{proposition}[theorem]{Proposition}
\newtheorem{lemma}[theorem]{Lemma}
\newtheorem{corollary}[theorem]{Corollary}
\theoremstyle{definition}
\newtheorem{definition}[theorem]{Definition}
\newtheorem{assumption}[theorem]{Assumption}
\theoremstyle{remark}
\newtheorem{remark}[theorem]{Remark}

\newcommand{\Sspace}{\mathcal S}
\newcommand{\Aspace}{\mathcal A}
\newcommand{\Xspace}{\mathcal X}
\newcommand{\Yspace}{\mathcal Y}
\newcommand{\Wspace}{\mathcal W}
\newcommand{\Pspace}{\mathcal P}
\newcommand{\Normal}{\mathcal N}
\newcommand{\Unif}{\operatorname{Unif}}
\newcommand{\TV}{\operatorname{TV}}
\newcommand{\Law}{\operatorname{Law}}
\newcommand{\E}{\mathbb E}
\newcommand{\R}{\mathbb R}
\newcommand{\nextx}{\operatorname{next}}
\newcommand{\stoptraj}{\operatorname{stop}}
\newcommand{\pushtraj}{\operatorname{push}}
\newcommand{\sg}{\operatorname{sg}}
\newcommand{\emptytok}{\bot}
\newcommand{\dd}{\mathrm d}
\newcommand{\score}{s}
\newcommand{\Bell}{\mathcal B}
\newcommand{\Diff}{\mathcal D}
\newcommand{\Student}{\theta}
\newcommand{\Teacher}{\bar\theta}
\newcommand{\targetpi}{\pi_e}
\newcommand{\branchstop}{\rightarrow}
\newcommand{\branchcont}{\curvearrowright}

\title{Bellman Path Diffusion for Direct Off-Policy Trajectory Generation}
\author{}
\date{}

\begin{document}
\maketitle

\begin{abstract}
We study off-policy evaluation from one-step transition data when the desired
output is a complete reward-bearing trajectory under a fixed evaluation
policy. We introduce Bellman Path Diffusion, a trajectory-space extension of
the probability-path construction in Temporal Difference Flows. A target-policy
path with maximum capacity $h$ is defined recursively: one real transition is
always generated, after which the path stops with probability $1-\gamma$ or
continues with a target-policy suffix with probability $\gamma$. Independent
blockwise corruption preserves this decomposition at every diffusion time and
gives both branches the same Gaussian source. This yields a single score
regression objective. The stopping target is analytic, while the continuation
target concatenates the analytic score of the current transition with the
frozen score of the future suffix at the same noise level. Under population
realizability, stagewise training recovers the exact noisy path law, and one
reverse diffusion directly generates the full padded trajectory tensor.
Geometric survival implements discounting, so the undiscounted reward sum of
the generated path is an unbiased estimator of the truncated discounted
return. We further prove full-law stability in total variation, a
return-sensitive contraction and OPE error bound, and an exact replay
equivalence that permits clean suffix caching without changing the population
noise-level objective.
\end{abstract}

\section{Introduction}
\label{sec:introduction}

Off-policy evaluation asks for the performance of a policy using data collected
by other policies. Classical estimators compress the future into a scalar value
or an occupancy correction. A generative alternative is to learn the complete
future distribution and evaluate arbitrary reward or safety functionals on
samples from that distribution. Standard world models are poorly matched to
this objective because they are unrolled one environment step at a time at
inference, causing model error to compound with horizon.

Geometric horizon models avoid environment-time unrolling by directly
sampling a discounted future state. Temporal Difference Flows
\citep{farebrother2025tdflow} makes this idea compatible with modern flow and
diffusion models. Its central observation is not merely that one can train a
generative model on Bellman-backed samples. The stronger observation is that
the two components of the successor-measure Bellman equation define two
probability paths with a common source. Regressing their vector fields or
scores with weights $1-\gamma$ and $\gamma$ directly constructs the
probability path of the Bellman mixture.

A future-state sample, however, is not a trajectory. It does not preserve the
joint dependence among rewards, states, and evaluation-policy actions, and it
cannot evaluate path-dependent quantities without additional assumptions. We
therefore lift the probability-path Bellman construction from successor states
to reward-bearing trajectories. This lift is not obtained by first generating
clean teacher trajectories and then applying ordinary diffusion training. The
Bellman decomposition is imposed directly at every diffusion time.

Our construction has one central identity. Let $m_{h,t}^{\targetpi}$ denote
the noisy law at diffusion time $t$ of an evaluation-policy trajectory with at
most $h$ reward-bearing transitions. Then
\begin{equation}
\boxed{
 m_{h,t}^{\targetpi}
 =
 (1-\gamma)\,m_{h,t}^{\branchstop}
 +
 \gamma\,m_{h,t}^{\branchcont},
 \qquad t\in[0,T].
}
\label{eq:intro-mixture}
\end{equation}
The stopping path contains one real transition followed by padding. The
continuation path contains one real transition followed by a noisy suffix drawn
from the $(h-1)$-step teacher at the same diffusion time. Both paths start from
the same Gaussian source at $t=T$. Their component scores are tractable, so
\cref{eq:intro-mixture} induces a single score-matching objective whose
population minimizer is $\nabla\log m_{h,t}^{\targetpi}$. At no point is a
sampling-time policy guidance term added to the reverse process.

Our contributions are as follows.
\begin{enumerate}
    \item We define a finite-depth target-policy path recursion whose
    nonpadding blocks form a valid reward-bearing trajectory and whose
    undiscounted reward sum represents the truncated discounted return.
    \item We prove that independent blockwise corruption commutes with this
    recursion. The stopping and continuation branches therefore form an exact
    mixture at every noise level and share a common Gaussian source.
    \item We derive the corresponding score backup. The current-transition
    score is analytic, the suffix score is supplied by a frozen shorter-horizon
    teacher, and their branch regression recovers the score of the target
    mixture without an auxiliary Bellman penalty or sampling-time guidance.
    \item We establish exact stagewise recovery, replay-buffer equivalence,
    full-law and return-sensitive error propagation, and direct OPE from one
    reverse diffusion.
\end{enumerate}

\begin{remark}[Terminology and scope]
\label{rem:terminology}
The word Bellman in this paper refers to the finite-depth path recursion that
prepends one real transition to a shorter future suffix. It is not the
stationary discounted-occupancy flow equation used by DICE methods. The core
technical statement is that this path recursion persists under blockwise
corruption and hence determines a score recursion at every diffusion time.
\end{remark}

\section{Problem Setup}
\label{sec:setup}

Consider a discounted Markov decision process
\begin{equation}
\mathcal M
=
(\Sspace,\Aspace,P^\star,\mu_0,\gamma),
\qquad
\gamma\in(0,1),
\label{eq:mdp}
\end{equation}
and a fixed evaluation policy $\targetpi(\dd a\mid s)$. Let
\begin{equation}
\Xspace:=\Sspace\times\Aspace,
\qquad
x=(s,a)\in\Xspace.
\end{equation}
The joint reward-transition kernel is
\begin{equation}
P^\star(\dd r,\dd s'\mid s,a).
\label{eq:environment-kernel}
\end{equation}
Define a reward-bearing transition token
\begin{equation}
y=(r,s',a')\in
\Yspace:=\R\times\Sspace\times\Aspace,
\end{equation}
and the evaluation-policy token kernel
\begin{equation}
\boxed{
\kappa_{\targetpi}(\dd y\mid x)
:=
P^\star(\dd r,\dd s'\mid s,a)
\targetpi(\dd a'\mid s').
}
\label{eq:token-kernel}
\end{equation}
We write
\begin{equation}
\nextx(y):=(s',a').
\end{equation}
The offline dataset is
\begin{equation}
\mathcal D
=
\{(s_i,a_i,r_i,s_i')\}_{i=1}^N
\end{equation}
and may have been collected by arbitrary behavior policies.

\begin{assumption}[Conditional validity and coverage]
\label{ass:coverage}
Conditioned on $X=(s,a)$, every logged pair $(R,S')$ follows
$P^\star(\cdot,\cdot\mid s,a)$. Every state-action pair reached by
$\targetpi$ before the chosen truncation horizon $H$ lies in the support of
the dataset state-action distribution.
\end{assumption}

Under \cref{ass:coverage}, drawing $(s,a,r,s')$ from the dataset and then
sampling
\begin{equation}
a'\sim\targetpi(\cdot\mid s')
\end{equation}
produces an unbiased sample $Y\sim\kappa_{\targetpi}(\cdot\mid X)$. No
behavior-policy density or importance weight is needed.

\section{Evaluation-Policy Bellman Recursion on Path Space}
\label{sec:path-flow}

\subsection{Stopped reward-bearing trajectories}

Let $\emptytok$ be a padding token. For $h\geq 1$, define
\begin{equation}
\Wspace_h
:=
(\Yspace\cup\{\emptytok\})^h,
\qquad
\Wspace_0:=\{\varnothing\}.
\end{equation}
A valid element contains a sequence of transition tokens followed by padding.
For $y\in\Yspace$, define
\begin{equation}
\stoptraj_h(y)
:=
(y,\emptytok,\ldots,\emptytok)
\in\Wspace_h.
\label{eq:stop-map}
\end{equation}
For $w=(w_0,\ldots,w_{h-2})\in\Wspace_{h-1}$, define
\begin{equation}
\pushtraj_h(y,w)
:=
(y,w_0,\ldots,w_{h-2})
\in\Wspace_h.
\label{eq:push-map}
\end{equation}
The empty case satisfies $\pushtraj_1(y,\varnothing)=(y)$.

Let $L_h(w)$ be the number of nonpadding tokens in $w$. Once padding occurs,
all later blocks are padding. Environment termination is represented by an
absorbing state-action pair with zero future rewards.

\subsection{Finite-depth Bellman recursion}

Define the conditional path laws recursively by
\begin{equation}
\mathbb M_0^{\targetpi}(\cdot\mid x)
:=
\delta_{\varnothing},
\label{eq:path-base}
\end{equation}
and, for $h\geq 1$,
\begin{equation}
\boxed{
\begin{aligned}
\mathbb M_h^{\targetpi}(\dd w\mid x)
:=
\int \kappa_{\targetpi}(\dd y\mid x)
\Big[
&(1-\gamma)
\delta_{\stoptraj_h(y)}(\dd w)
\\
&+
\gamma
(\pushtraj_h(y,\cdot))_{\#}
\mathbb M_{h-1}^{\targetpi}
(\dd w\mid\nextx(y))
\Big].
\end{aligned}
}
\label{eq:path-bellman}
\end{equation}
The stopping branch injects one real evaluation-policy transition. The
continuation branch prepends the same transition to a future suffix. We denote
the operator mapping a conditional law on $\Wspace_{h-1}$ to one on
$\Wspace_h$ by $\Bell_h^{\targetpi}$, so that
\begin{equation}
\mathbb M_h^{\targetpi}
=
\Bell_h^{\targetpi}\mathbb M_{h-1}^{\targetpi}.
\end{equation}

\begin{proposition}[Trajectory semantics]
\label{prop:path-semantics}
Let $W_h\sim\mathbb M_h^{\targetpi}(\cdot\mid X_0=x_0)$. Conditional on
$L_h(W_h)>t$, its valid tokens satisfy
\begin{equation}
(R_t,S_{t+1})
\sim
P^\star(\cdot,\cdot\mid S_t,A_t),
\qquad
A_{t+1}
\sim
\targetpi(\cdot\mid S_{t+1}).
\label{eq:valid-path}
\end{equation}
Moreover,
\begin{equation}
\Pr(L_h(W_h)>t\mid X_0=x_0)=\gamma^t,
\qquad
0\leq t\leq h-1.
\label{eq:survival}
\end{equation}
Equivalently, the censored length distribution is
\begin{equation}
\boxed{
\Pr(L_h=\ell)
=
\begin{cases}
(1-\gamma)\gamma^{\ell-1}, & 1\leq \ell<h,\\
\gamma^{h-1}, & \ell=h.
\end{cases}
}
\label{eq:length-law}
\end{equation}
Consequently,
\begin{equation}
\boxed{
\E_{W_h\sim\mathbb M_h^{\targetpi}(\cdot\mid x_0)}
\left[
\sum_{t=0}^{L_h(W_h)-1}R_t
\right]
=
\E_{\targetpi}
\left[
\sum_{t=0}^{h-1}\gamma^tR_t
\middle|
X_0=x_0
\right].
}
\label{eq:value-identity}
\end{equation}
\end{proposition}

\begin{proof}
Equation~\eqref{eq:valid-path} follows by recursively expanding
\cref{eq:path-bellman}. A token at depth $t$ exists exactly when the first
$t$ continuation decisions succeed, which has probability $\gamma^t$. Hence
\begin{align}
\E\left[\sum_{t=0}^{L_h(W_h)-1}R_t\right]
&=
\sum_{t=0}^{h-1}
\E\left[\bm 1\{L_h(W_h)>t\}R_t\right]
\\
&=
\sum_{t=0}^{h-1}
\gamma^t\E_{\targetpi}[R_t\mid X_0=x_0].
\end{align}
\end{proof}

\begin{remark}[Why geometric stopping is useful]
The random stopping rule is the probabilistic representation of
discounting used by normalized successor measures. It converts a discounted
sum into an undiscounted reward sum on a random-length path. The fixed tensor
capacity $h$ only censors paths that survive beyond the computational horizon;
the last length bin in \cref{eq:length-law} therefore contains all surviving
mass.
\end{remark}

\section{Bellman Probability Paths}
\label{sec:probability-paths}

\subsection{Blockwise forward diffusion}

Let
\begin{equation}
\varphi:
\Yspace\cup\{\emptytok\}
\to
\R^d
\end{equation}
be an injective token representation. For
$w=(w_0,\ldots,w_{h-1})$, write
\begin{equation}
\Phi_h(w)
:=
(\varphi(w_0),\ldots,\varphi(w_{h-1}))
\in\R^{hd}.
\end{equation}
Consider a linear forward diffusion with one-block kernel
\begin{equation}
q_{t\mid0}(\dd z\mid u)
=
\Normal
\left(
\dd z;
\alpha_t\varphi(u),
\sigma_t^2I_d
\right),
\qquad
t\in[0,T].
\label{eq:block-noise}
\end{equation}
We assume
\begin{equation}
\alpha_0=1,
\quad
\sigma_0=0,
\quad
\alpha_T=0,
\quad
\sigma_T=1.
\label{eq:boundary-schedule}
\end{equation}
The $h$-block corruption kernel factorizes as
\begin{equation}
\boxed{
q_{t\mid0}^{(h)}(\dd z_{0:h-1}\mid w)
:=
\prod_{j=0}^{h-1}
q_{t\mid0}(\dd z_j\mid w_j).
}
\label{eq:blockwise-kernel}
\end{equation}
Let
\begin{equation}
m_{h,t}^{\targetpi}(\cdot\mid x)
:=
(\Diff_{h,t})_{\#}
\mathbb M_h^{\targetpi}(\cdot\mid x)
\label{eq:noisy-path-law}
\end{equation}
be the resulting noisy trajectory law.

\subsection{The two Bellman branches at every noise level}

Define the stopping probability path
\begin{equation}
\boxed{
\begin{aligned}
m_{h,t}^{\branchstop}(\dd z\mid x)
:=
\int
\kappa_{\targetpi}(\dd y\mid x)
q_{t\mid0}^{(h)}
(\dd z\mid\stoptraj_h(y)).
\end{aligned}
}
\label{eq:stop-path}
\end{equation}
Write $z=(z_0,z_+)\in\R^d\times\R^{(h-1)d}$. Define the continuation
probability path
\begin{equation}
\boxed{
\begin{aligned}
m_{h,t}^{\branchcont}
(\dd z_0,\dd z_+\mid x)
:=
\int
\kappa_{\targetpi}(\dd y\mid x)
q_{t\mid0}(\dd z_0\mid y)
 m_{h-1,t}^{\targetpi}
(\dd z_+\mid\nextx(y)).
\end{aligned}
}
\label{eq:continue-path}
\end{equation}
For $h=1$, the suffix space is a singleton and the two paths coincide.

\begin{theorem}[Trajectory Bellman probability path]
\label{thm:path-mixture}
For every $h\geq 1$, every conditioning pair $x$, and every diffusion time
$t\in[0,T]$,
\begin{equation}
\boxed{
 m_{h,t}^{\targetpi}(\cdot\mid x)
 =
 (1-\gamma)
 m_{h,t}^{\branchstop}(\cdot\mid x)
 +
 \gamma
 m_{h,t}^{\branchcont}(\cdot\mid x).
}
\label{eq:noise-bellman}
\end{equation}
At the clean endpoint,
\begin{equation}
m_{h,0}^{\targetpi}
=
(\Phi_h)_{\#}\mathbb M_h^{\targetpi}.
\label{eq:clean-endpoint}
\end{equation}
At the source endpoint,
\begin{equation}
m_{h,T}^{\branchstop}
=
m_{h,T}^{\branchcont}
=
m_{h,T}^{\targetpi}
=
\Normal(0,I_{hd}).
\label{eq:common-source}
\end{equation}
\end{theorem}

\begin{proof}
Substitute the path recursion in \cref{eq:path-bellman} into the pushforward
in \cref{eq:noisy-path-law}. On the stopping branch this gives
\cref{eq:stop-path}. On the continuation branch, blockwise corruption makes
the noisy current token independent of the noisy suffix conditioned on the
clean path. Integrating over the clean suffix gives
$m_{h-1,t}^{\targetpi}$, which yields \cref{eq:continue-path}. At $t=T$,
\cref{eq:boundary-schedule} makes every block standard Gaussian independently
of its clean token. Both branches therefore share the same source.
\end{proof}

\begin{remark}[The source-boundary issue]
A clean Bellman equation should not be copied mechanically to every noise
level. The correct object is the time-dependent decomposition in
\cref{eq:noise-bellman}. Each branch is transported through its own diffusion
probability path, and the paths meet at the common Gaussian source. This is the
trajectory analogue of the probability-path Bellman construction in Temporal
Difference Flows.
\end{remark}

\section{Direct Bellmanization of the Reverse Score}
\label{sec:score-learning}

The reverse SDE associated with the forward process has drift
\begin{equation}
\dd Z_t
=
\left[
 f(t)Z_t
 -
 g(t)^2\nabla_z\log m_{h,t}^{\targetpi}(Z_t\mid x)
\right]\dd t
+
g(t)\dd\overline W_t,
\qquad
t:T\downarrow0.
\label{eq:reverse-sde}
\end{equation}
Thus direct trajectory generation reduces to learning the score of
\cref{eq:noise-bellman}.

\subsection{A mixture-score identity}

\begin{lemma}[Score of a probability-path mixture]
\label{lem:mixture-score}
Let $p^0$ and $p^1$ be differentiable densities and
$p=(1-\gamma)p^0+\gamma p^1$. Draw a branch
$B\sim\operatorname{Bernoulli}(\gamma)$ and then draw $Z\sim p^B$. If
$U^b(z)=\nabla_z\log p^b(z)$, then
\begin{equation}
\boxed{
\arg\min_{s}
\E\left[\|s(Z)-U^B(Z)\|_2^2\right]
=
\nabla_z\log p(z).
}
\label{eq:mixture-score-regression}
\end{equation}
The minimization is over all square-integrable vector fields.
\end{lemma}

\begin{proof}
The population minimizer of squared regression is the conditional mean
$\E[U^B(Z)\mid Z=z]$. Bayes' rule gives
\begin{align}
\E[U^B(Z)\mid Z=z]
&=
\frac{
(1-\gamma)p^0(z)\nabla\log p^0(z)
+\gamma p^1(z)\nabla\log p^1(z)
}{p(z)}
\\
&=
\frac{\nabla p(z)}{p(z)}
=
\nabla\log p(z).
\end{align}
\end{proof}

\subsection{Tractable component scores}

For the stopping branch, sample
\begin{equation}
Y\sim\kappa_{\targetpi}(\cdot\mid X),
\qquad
Z^{\branchstop}_{0:h-1}
\sim
q_{t\mid0}^{(h)}
(\cdot\mid\stoptraj_h(Y)).
\label{eq:stop-sampling}
\end{equation}
Denoising score matching permits the conditional target
\begin{equation}
\boxed{
U_{h,t}^{\branchstop}
:=
\nabla_z
\log
q_{t\mid0}^{(h)}
\left(
Z^{\branchstop}_{0:h-1}
\mid
\stoptraj_h(Y)
\right).
}
\label{eq:stop-score-target}
\end{equation}
For the Gaussian kernel,
\begin{equation}
U_{h,t}^{\branchstop}
=
-
\frac{1}{\sigma_t^2}
\begin{pmatrix}
Z_0^{\branchstop}-\alpha_t\varphi(Y)
\\
Z_1^{\branchstop}-\alpha_t\varphi(\emptytok)
\\
\vdots
\\
Z_{h-1}^{\branchstop}-\alpha_t\varphi(\emptytok)
\end{pmatrix}.
\label{eq:stop-gaussian-score}
\end{equation}

For the continuation branch, draw the same real token
$Y\sim\kappa_{\targetpi}(\cdot\mid X)$ and set
$X'=\nextx(Y)$. Draw
\begin{equation}
Z_0^{\branchcont}
\sim
q_{t\mid0}(\cdot\mid Y).
\label{eq:cont-current-noise}
\end{equation}
The suffix is sampled directly from the frozen teacher probability path at the
same noise level:
\begin{equation}
Z_+^{\branchcont}
\sim
m_{h-1,t}^{\Teacher}
(\cdot\mid X').
\label{eq:partial-reverse-suffix}
\end{equation}
Operationally, \cref{eq:partial-reverse-suffix} starts from Gaussian noise at
$T$ and integrates the teacher reverse SDE only until $t$. It does not require
generating a clean suffix and corrupting it again.

Let
\begin{equation}
\score_{\Teacher,h-1}
(z_+,t\mid x')
\approx
\nabla_{z_+}
\log
m_{h-1,t}^{\targetpi}(z_+\mid x')
\end{equation}
be the frozen suffix score. The continuation target is
\begin{equation}
\boxed{
U_{h,t}^{\branchcont}
:=
\begin{pmatrix}
\nabla_{z_0}
\log q_{t\mid0}
(Z_0^{\branchcont}\mid Y)
\\[1mm]
\sg\!
\left[
\score_{\Teacher,h-1}
(Z_+^{\branchcont},t\mid X')
\right]
\end{pmatrix}.
}
\label{eq:continuation-score-target}
\end{equation}
The first block is analytic. The remaining blocks are bootstrapped at the same
diffusion time from the previous horizon.

\subsection{The single Bellman diffusion objective}

Let
\begin{equation}
\score_{\Student,h}
(z,t\mid x)
\end{equation}
be a trajectory score model. One transformer may be shared across all $h$ by
using a horizon token and a variable-length attention mask. Define
\begin{equation}
\boxed{
\begin{aligned}
\mathcal L_h(\Student;\Teacher)
:=
\E_{X,t}
\Big[
&(1-\gamma)
\E
\left[
\lambda(t)
\left\|
\score_{\Student,h}
(Z^{\branchstop},t\mid X)
-
U_{h,t}^{\branchstop}
\right\|_2^2
\right]
\\
&+
\gamma
\E
\left[
\lambda(t)
\left\|
\score_{\Student,h}
(Z^{\branchcont},t\mid X)
-
U_{h,t}^{\branchcont}
\right\|_2^2
\right]
\Big].
\end{aligned}
}
\label{eq:main-loss}
\end{equation}
Here $X$ is sampled from the dataset state-action distribution and
$t\sim p_T$ is sampled from the diffusion training schedule. In practice, one
samples a branch indicator and evaluates only its corresponding term.

\begin{remark}[There is no auxiliary base loss]
Equation~\eqref{eq:main-loss} is not a regularizer added to an ordinary
diffusion objective. It is the denoising score-matching objective of the
Bellman mixture itself. Removing either the stopping or continuation term
changes the target probability path and therefore changes the generated clean
law.
\end{remark}

\begin{theorem}[Population score backup]
\label{thm:population-score-backup}
Assume the frozen teacher is exact at horizon $h-1$:
\begin{equation}
\score_{\Teacher,h-1}(z,t\mid x)
=
\nabla_z\log m_{h-1,t}^{\targetpi}(z\mid x).
\label{eq:exact-teacher}
\end{equation}
Then every population minimizer of \cref{eq:main-loss} satisfies
\begin{equation}
\boxed{
\score_{\Student,h}^\star(z,t\mid x)
=
\nabla_z\log m_{h,t}^{\targetpi}(z\mid x)
}
\label{eq:exact-student-score}
\end{equation}
for almost every $(z,t,x)$.
\end{theorem}

\begin{proof}
For the stopping branch, the denoising score-matching identity converts the
conditional target in \cref{eq:stop-score-target} into the marginal score
$\nabla\log m_{h,t}^{\branchstop}$. For the continuation branch, conditional
on $(Y,X')$, the joint density factorizes as
\begin{equation}
q_{t\mid0}(z_0\mid Y)
 m_{h-1,t}^{\targetpi}(z_+\mid X').
\end{equation}
Its score is exactly \cref{eq:continuation-score-target}. Marginalizing
$(Y,X')$ therefore gives $\nabla\log m_{h,t}^{\branchcont}$. Applying
\cref{lem:mixture-score} with the weights in \cref{eq:noise-bellman} yields
\cref{eq:exact-student-score}.
\end{proof}

\section{Replay-Buffered Noisy Suffix Sampling}
\label{sec:replay}

The continuation law in \cref{eq:continue-path} requires a suffix sample from
the teacher marginal at the same diffusion time. There are two exact
population implementations when the teacher score and reverse solver are
exact. One may partially integrate the teacher reverse SDE from $T$ to $t$, or
one may generate a clean teacher suffix once, cache it, and independently
corrupt it at each requested noise level.

Let
\begin{equation}
\widehat{\mathbb M}_{\Teacher,h-1}(\cdot\mid x')
\end{equation}
denote the clean suffix law induced by a frozen teacher, and define its
forward-corrupted marginal by
\begin{equation}
\widehat m_{\Teacher,h-1,t}(\cdot\mid x')
:=
(\Diff_{h-1,t})_{\#}
\widehat{\mathbb M}_{\Teacher,h-1}(\cdot\mid x').
\label{eq:teacher-forward-marginal}
\end{equation}
A replay entry stores the conditioning pair and the clean suffix:
\begin{equation}
(x',W_+),
\qquad
W_+\sim
\widehat{\mathbb M}_{\Teacher,h-1}(\cdot\mid x').
\label{eq:cached-suffix}
\end{equation}

\begin{proposition}[Replay equivalence]
\label{prop:replay-equivalence}
Conditioned on the same successor pair $x'$, draw a cached clean suffix as in
\cref{eq:cached-suffix} and independently apply the blockwise corruption
kernel at time $t$. Then
\begin{equation}
\boxed{
Z_+
\sim
\widehat m_{\Teacher,h-1,t}(\cdot\mid x').
}
\label{eq:replay-equivalence}
\end{equation}
Consequently, clean-suffix replay gives exactly the same population
continuation objective as sampling from the teacher forward probability path.
If the teacher is an exact score model for that path and its reverse SDE is
solved exactly, partial reverse integration to time $t$ has the same marginal
and is equivalent as well.
\end{proposition}

\begin{proof}
Equation~\eqref{eq:teacher-forward-marginal} is the pushforward of the clean
teacher law through the independent blockwise corruption kernel. Sampling a
clean suffix and applying that kernel is therefore sampling from the stated
pushforward. Under score consistency and exact reverse integration, the
teacher reverse process has these same marginals.
\end{proof}

\begin{remark}[What caching does and does not claim]
\label{rem:caching-claim}
The theory does not require avoiding clean suffix generation. Its substantive
claim is that the Bellman path decomposition and score backup are defined
exactly at every noise level. Replay is a computational implementation of that
noise-level law. For continuous conditioning variables, a cached suffix must
remain attached to the same sampled $x'=(s',a')$; nearest-neighbor reuse changes
the conditional target and is an additional approximation.
\end{remark}

\section{Exact Recovery and Error Propagation}
\label{sec:theory}

\begin{assumption}[Population realizability]
\label{ass:realizability}
For every horizon $h\leq H$, the score model class contains
$\nabla\log m_{h,t}^{\targetpi}(\cdot\mid x)$ for almost every $(t,x)$, the
objective is minimized exactly, and the reverse SDE is solved exactly.
\end{assumption}

\begin{theorem}[Stagewise exact recovery]
\label{thm:stagewise-recovery}
Under \cref{ass:coverage,ass:realizability}, train horizons in the order
$h=1,\ldots,H$, using the exact horizon-$(h-1)$ score as the frozen teacher.
Then
\begin{equation}
\score_h(z,t\mid x)
=
\nabla_z\log m_{h,t}^{\targetpi}(z\mid x)
\end{equation}
for every $h\leq H$. Consequently, solving the horizon-$H$ reverse SDE from
$\Normal(0,I_{Hd})$ returns
\begin{equation}
\boxed{
W_H
\sim
\mathbb M_H^{\targetpi}(\cdot\mid X_0).
}
\label{eq:direct-generation-law}
\end{equation}
\end{theorem}

\begin{proof}
For $h=1$, the suffix is empty, and the stopping and continuation branches
coincide with the one-token distribution. Its score is recovered by analytic
denoising score matching. Suppose the claim holds for $h-1$. Theorem
\ref{thm:population-score-backup} then gives the exact score at horizon $h$.
Induction proves the score claim for all $h\leq H$. The reverse-time SDE theorem
for score-based generative models then transports the common Gaussian source
to the clean endpoint $(\Phi_H)_{\#}\mathbb M_H^{\targetpi}$.
\end{proof}

The stagewise construction differs from a generic fitted fixed-point iteration.
The teacher depth is the trajectory suffix length. Each stage adds exactly one
real evaluation-policy transition in front of the previously learned suffix.
At inference, however, the full $H$-step tensor is produced by one reverse
process.

\subsection{Approximate stagewise fitting}

For conditional laws on $\Wspace_h$, define
\begin{equation}
\|Q-Q'\|_{\TV,\infty}
:=
\sup_{x\in\Xspace}
\TV(Q(\cdot\mid x),Q'(\cdot\mid x)).
\end{equation}

\begin{lemma}[Full-law stability]
\label{lem:bellman-stability}
For conditional suffix laws $Q,Q'$ on $\Wspace_{h-1}$,
\begin{equation}
\boxed{
\|\Bell_h^{\targetpi}Q-\Bell_h^{\targetpi}Q'\|_{\TV,\infty}
\leq
\gamma
\|Q-Q'\|_{\TV,\infty}.
}
\label{eq:bellman-tv-stability}
\end{equation}
\end{lemma}

\begin{proof}
The stopping branches are identical. The continuation branches differ only
through the suffix laws, have total mass $\gamma$, and are obtained by
integration and measurable pushforward, both of which are non-expansive in
total variation.
\end{proof}

Let $\widehat{\mathbb M}_h$ be the clean law induced by the learned
horizon-$h$ reverse process, with
$\widehat{\mathbb M}_0=\mathbb M_0^{\targetpi}$. Define
\begin{equation}
\varepsilon_h^{\TV}
:=
\left\|
\widehat{\mathbb M}_h
-
\Bell_h^{\targetpi}\widehat{\mathbb M}_{h-1}
\right\|_{\TV,\infty}.
\label{eq:tv-stage-error}
\end{equation}

\begin{theorem}[Full trajectory-law error]
\label{thm:tv-error}
For every $h\leq H$,
\begin{equation}
\boxed{
\left\|
\widehat{\mathbb M}_h-\mathbb M_h^{\targetpi}
\right\|_{\TV,\infty}
\leq
\sum_{j=1}^{h}\gamma^{h-j}\varepsilon_j^{\TV}.
}
\label{eq:tv-error-bound}
\end{equation}
\end{theorem}

\begin{proof}
Apply the triangle inequality, \cref{lem:bellman-stability}, and the exact
recursion
$\mathbb M_h^{\targetpi}
=\Bell_h^{\targetpi}\mathbb M_{h-1}^{\targetpi}$:
\[
\delta_h^{\TV}
\leq
\varepsilon_h^{\TV}
+
\gamma\delta_{h-1}^{\TV}.
\]
Unrolling from $\delta_0^{\TV}=0$ gives
\cref{eq:tv-error-bound}.
\end{proof}

\subsection{Return-sensitive seminorm}

Full-law error is stronger than OPE requires. Define the undiscounted reward
sum on a stopped path by
\begin{equation}
G_h(w)
:=
\sum_{t=0}^{L_h(w)-1}r(w_t),
\qquad
G_0(\varnothing):=0,
\label{eq:return-functional}
\end{equation}
where $r(y)$ denotes the reward coordinate of token $y$. For a conditional law
$Q$ on $\Wspace_h$, define
\begin{equation}
V_h^Q(x)
:=
\int G_h(w)Q(\dd w\mid x).
\label{eq:law-value}
\end{equation}
For two conditional laws, define
\begin{equation}
\boxed{
\|Q-Q'\|_{G,h}
:=
\sup_{x\in\Xspace}
\left|
V_h^Q(x)-V_h^{Q'}(x)
\right|.
}
\label{eq:g-seminorm}
\end{equation}
This is a seminorm because distinct path laws may have the same expected
reward sum.

The path recursion implies
\begin{equation}
\boxed{
V_h^{\Bell_h^{\targetpi}Q}(x)
=
\int\kappa_{\targetpi}(\dd y\mid x)
\left[
r(y)
+
\gamma V_{h-1}^{Q}(\nextx(y))
\right].
}
\label{eq:value-bellman-recursion}
\end{equation}

\begin{theorem}[Return-sensitive contraction]
\label{thm:g-contraction}
For any two conditional suffix laws $Q,Q'$,
\begin{equation}
\boxed{
\|\Bell_h^{\targetpi}Q-\Bell_h^{\targetpi}Q'\|_{G,h}
\leq
\gamma\|Q-Q'\|_{G,h-1}.
}
\label{eq:g-contraction}
\end{equation}
\end{theorem}

\begin{proof}
The immediate reward in \cref{eq:value-bellman-recursion} is identical for
$Q$ and $Q'$ and cancels. Taking absolute values and the supremum over $x$
leaves $\gamma\|Q-Q'\|_{G,h-1}$.
\end{proof}

Define the value-sensitive stage error
\begin{equation}
\varepsilon_h^G
:=
\left\|
\widehat{\mathbb M}_h
-
\Bell_h^{\targetpi}\widehat{\mathbb M}_{h-1}
\right\|_{G,h}.
\label{eq:g-stage-error}
\end{equation}
Equivalently,
\begin{equation}
\boxed{
\varepsilon_h^G
=
\sup_x
\left|
\widehat V_h(x)
-
\int\kappa_{\targetpi}(\dd y\mid x)
\left[
r(y)+\gamma\widehat V_{h-1}(\nextx(y))
\right]
\right|,
}
\label{eq:g-bellman-residual}
\end{equation}
where $\widehat V_h:=V_h^{\widehat{\mathbb M}_h}$.

\begin{theorem}[Return error propagation]
\label{thm:g-error}
Let
\begin{equation}
\delta_h^G
:=
\left\|
\widehat{\mathbb M}_h-\mathbb M_h^{\targetpi}
\right\|_{G,h}.
\end{equation}
Then
\begin{equation}
\boxed{
\delta_h^G
\leq
\varepsilon_h^G+\gamma\delta_{h-1}^G,
\qquad
\delta_H^G
\leq
\sum_{h=1}^{H}\gamma^{H-h}\varepsilon_h^G.
}
\label{eq:g-error-bound}
\end{equation}
If $\varepsilon_h^G\leq\varepsilon^G$ for all $h$, then
\begin{equation}
\boxed{
\delta_H^G
\leq
\frac{1-\gamma^H}{1-\gamma}\varepsilon^G.
}
\label{eq:g-uniform-bound}
\end{equation}
\end{theorem}

\begin{proof}
Use the triangle inequality, the exact path recursion, and
\cref{thm:g-contraction}. This gives the one-step recursion in
\cref{eq:g-error-bound}; repeated substitution and the finite geometric series
give the remaining statements.
\end{proof}

\begin{remark}[Two levels of control]
\label{rem:two-errors}
Total variation controls the complete generated path law. The $G$-seminorm
only controls conditional expected returns, but it gives the OPE quantity
directly and avoids multiplying the final error by the maximum tensor length.
Which error decreases or grows with horizon in finite models is an empirical
question; \cref{thm:g-error} isolates the stagewise quantities that should be
measured.
\end{remark}

\section{Training and Generation}
\label{sec:algorithm}

\subsection{Stagewise training}

The exact theory uses one frozen teacher for each shorter suffix depth. The
practical version may obtain the teacher noisy suffix either by partial reverse
integration or by the replay equivalence in \cref{prop:replay-equivalence}.
The latter amortizes expensive full teacher generations over many score
updates.

\begin{algorithm}[H]
\caption{Bellman Path Diffusion Training with Suffix Replay}
\label{alg:training}
\begin{algorithmic}[1]
\Require Offline transitions $\mathcal D$, evaluation policy $\targetpi$,
discount $\gamma$, maximum capacity $H$, diffusion schedule, score network
$\score_\theta(z,t\mid x,h)$
\Ensure Horizon-$H$ trajectory score
\State Define the horizon-$0$ suffix as the empty law
\For{$h=1,\ldots,H$}
    \State Freeze a target score $\score_{\Teacher,h-1}$ and its clean
    horizon-$(h-1)$ reverse generator
    \State Initialize a replay buffer $\mathcal R_{h-1}$ keyed by the exact
    conditioning pair $x'$
    \State Refresh $\mathcal R_{h-1}$ by sampling logged successors,
    target actions, and clean teacher suffixes
    \Repeat
        \State Sample $(s,a,r,s')\sim\mathcal D$ and set $x\gets(s,a)$
        \State Sample $a'\sim\targetpi(\cdot\mid s')$,
        $x'\gets(s',a')$, and $y\gets(r,s',a')$
        \State Sample a diffusion time $t$ and
        $C\sim\operatorname{Bernoulli}(\gamma)$
        \If{$C=0$ or $h=1$}
            \State Set $w\gets\stoptraj_h(y)$
            \State Sample $z\sim q_{t\mid0}^{(h)}(\cdot\mid w)$
            \State Set
            $u\gets\nabla_z\log q_{t\mid0}^{(h)}(z\mid w)$
        \Else
            \State Retrieve or generate
            $w_+\sim
            \widehat{\mathbb M}_{\Teacher,h-1}(\cdot\mid x')$
            and keep it attached to this $x'$
            \State Sample
            $z_0\sim q_{t\mid0}(\cdot\mid y)$ and
            $z_+\sim q_{t\mid0}^{(h-1)}(\cdot\mid w_+)$
            \State Set
            \[
            z\gets(z_0,z_+),
            \qquad
            u\gets
            \begin{pmatrix}
            \nabla_{z_0}\log q_{t\mid0}(z_0\mid y)
            \\
            \sg[\score_{\Teacher,h-1}(z_+,t\mid x')]
            \end{pmatrix}
            \]
        \EndIf
        \State Update $\theta$ using
        $\lambda(t)\|\score_\theta(z,t\mid x,h)-u\|_2^2$
    \Until{the horizon-$h$ training budget is exhausted}
    \State Freeze or EMA-copy the horizon-$h$ score for the next stage
\EndFor
\State \Return $\score_\theta(\cdot,\cdot\mid\cdot,H)$
\end{algorithmic}
\end{algorithm}

\begin{remark}[Shared network implementation]
The mathematical description uses a score family indexed by $h$. A single
variable-length transformer may share parameters across depths using a horizon
embedding and an attention mask. The stagewise theorem is the clean population
reference point; EMA targets and mixed-depth minibatches are two-timescale
approximations.
\end{remark}

\subsection{Direct generation and OPE}

After training, draw
\begin{equation}
S_0\sim\mu_0,
\qquad
A_0\sim\targetpi(\cdot\mid S_0),
\qquad
X_0=(S_0,A_0).
\end{equation}
Initialize
\begin{equation}
Z_T\sim\Normal(0,I_{Hd})
\end{equation}
and solve the horizon-$H$ reverse SDE in \cref{eq:reverse-sde}. Decode the
clean endpoint into $W_H\in\Wspace_H$ and remove all blocks after the first
padding token.

\begin{algorithm}[H]
\caption{Direct Evaluation-Policy Trajectory Generation}
\label{alg:generation}
\begin{algorithmic}[1]
\Require Initial distribution $\mu_0$, evaluation policy $\targetpi$,
trained horizon-$H$ score, reverse solver
\State Sample $s_0\sim\mu_0$ and $a_0\sim\targetpi(\cdot\mid s_0)$
\State Sample $z_T\sim\Normal(0,I_{Hd})$
\State Solve the reverse SDE from $T$ to $0$ conditioned on
$x_0=(s_0,a_0)$
\State Decode $w_H\gets\Phi_H^{-1}(z_0)$
\State Remove all blocks after the first padding token
\State \Return the generated reward-bearing trajectory
\end{algorithmic}
\end{algorithm}

Generate $M$ independent trajectories and estimate
\begin{equation}
\boxed{
\widehat J_H(\targetpi)
:=
\frac{1}{M}
\sum_{m=1}^M
\sum_{t=0}^{L_H(W_H^{(m)})-1}
R_t^{(m)}.
}
\label{eq:ope-estimator}
\end{equation}
No explicit factor $\gamma^t$ is required in the estimator because geometric
survival already implements discounting.

\begin{corollary}[Exact and approximate OPE]
\label{cor:ope}
Under the assumptions of \cref{thm:stagewise-recovery},
\begin{equation}
\E[\widehat J_H(\targetpi)]
=
J_H(\targetpi)
:=
\E_{\targetpi}
\left[
\sum_{t=0}^{H-1}\gamma^tR_t
\right].
\label{eq:exact-ope}
\end{equation}
More generally, let
\begin{equation}
\overline J_H(\targetpi)
:=
\E_{\substack{
X_0\sim\mu_0\targetpi\\
W_H\sim\widehat{\mathbb M}_H(\cdot\mid X_0)
}}
[G_H(W_H)].
\end{equation}
Then
\begin{equation}
\boxed{
|\overline J_H(\targetpi)-J_H(\targetpi)|
\leq
\delta_H^G
\leq
\sum_{h=1}^{H}\gamma^{H-h}\varepsilon_h^G.
}
\label{eq:approx-ope}
\end{equation}
If $|R_t|\leq R_{\max}$ and
$J(\targetpi)=\E_{\targetpi}[\sum_{t=0}^{\infty}\gamma^tR_t]$, then
\begin{equation}
\boxed{
|\overline J_H(\targetpi)-J(\targetpi)|
\leq
\sum_{h=1}^{H}\gamma^{H-h}\varepsilon_h^G
+
\frac{\gamma^H R_{\max}}{1-\gamma}.
}
\label{eq:infinite-ope}
\end{equation}
\end{corollary}

\begin{proof}
Exactness follows from \cref{prop:path-semantics,thm:stagewise-recovery}.
The first approximate bound is the definition of the $G$-seminorm averaged
over the initial state-action law, followed by \cref{thm:g-error}. The
infinite-horizon result adds the usual truncation tail.
\end{proof}

\section{Relationship to Temporal Difference Flows}
\label{sec:relationship}

The proposed method takes one specific technical idea from Temporal Difference
Flows: a Bellman decomposition can be lifted from the clean endpoint to the
complete generative probability path, and component vector fields or scores can
be bootstrapped directly. In the diffusion extension of
\citet{farebrother2025tdflow}, TD$^2$-DD uses an analytic score for the
one-step branch and the frozen score of the bootstrapped successor-state path
for the continuation branch. Our method preserves this construction exactly,
but changes the generated object and the continuation geometry.

The source method generates one geometrically sampled future state. We generate
a complete sequence of transition tokens. The continuation branch therefore
factorizes into
\begin{equation}
\boxed{
\text{analytic noisy current token}
\quad\oplus\quad
\text{teacher noisy future suffix}.
}
\end{equation}
This factorization is what makes trajectory generation possible without clean
pseudo-trajectories. It also produces a finite-depth induction: horizon $h$
uses horizon $h-1$ as its teacher, while a single horizon-$H$ reverse process
is used at inference.

The posterior-mean temporal-difference objective of
\citet{ying2026tddiffusion} is orthogonal. It improves consistency between two
noise levels for a known clean data distribution. It is not needed to define
the Bellman target distribution here, so we omit it from the core objective.
It may be added later as a solver-consistency regularizer, but it should not be
part of the correctness argument.

\section{Discussion}
\label{sec:discussion}

\paragraph{Why the framework is not a plug-in correction.}
The evaluation policy enters the one-step token kernel in
\cref{eq:token-kernel}. The Bellman stop and continuation branches define the
training distribution at every noise level through \cref{eq:noise-bellman}.
The resulting score is the field used by the ordinary reverse SDE. There is no
classifier guidance, density-ratio multiplier, behavior-policy correction, or
Bellman residual added during sampling.

\paragraph{Why complete trajectories matter.}
A successor-state model can evaluate rewards that are functions of one future
state. A reward-bearing path model preserves transition rewards, action
sequences, termination, cumulative costs, hitting events, and other path
functionals. Direct generation also avoids repeated calls to a learned
one-step world model at inference.

\paragraph{Frozen methodological scope.}
The core method consists only of the path recursion, blockwise corruption, the
two branch score targets, and stagewise suffix bootstrapping. Occupancy-ratio
estimation, Doob transforms, Schr\"odinger bridges, and sampling-time policy
guidance are outside this framework.

\paragraph{Main limitations.}
The method requires conditional support for state-action pairs reached by the
evaluation policy. Partial reverse teacher sampling adds training cost, and the
trajectory dimension grows linearly with $H$. Exact discrete token decoding is
also nontrivial when states or actions are categorical. These are
implementation and approximation challenges; they do not alter the
probability-path identity.

\paragraph{Empirical priorities.}
The first diagnostic is the generated length law in \cref{eq:length-law},
including the censored mass at $H$. The second is blockwise trajectory
validity under $P^\star$ and $\targetpi$. The third measures
$\varepsilon_h^{\TV}$ and $\varepsilon_h^G$ as functions of $h$ to determine
whether practical suffix errors are attenuated or accumulated. Replay and
partial reverse sampling should also be compared for computational cost and
finite-model mismatch, even though they are population-equivalent under the
conditions of \cref{prop:replay-equivalence}.

\section{Complete Operator Chain}
\label{sec:summary}

The complete construction is
\begin{equation}
\boxed{
\begin{aligned}
&Y=(R,S',A'),
\qquad
Y\sim P^\star(\dd r,\dd s'\mid S,A)
\targetpi(\dd a'\mid S')
\\[1mm]
&\Downarrow
\\[-1mm]
&\mathbb M_h^{\targetpi}
=
\Bell_h^{\targetpi}\mathbb M_{h-1}^{\targetpi}
\\[1mm]
&\Downarrow
\\[-1mm]
&m_{h,t}^{\targetpi}
=
(1-\gamma)m_{h,t}^{\branchstop}
+
\gamma m_{h,t}^{\branchcont}
\qquad\forall t
\\[1mm]
&\Downarrow
\\[-1mm]
&U_{h,t}^{\branchstop}
=
\text{analytic full-path score},
\qquad
U_{h,t}^{\branchcont}
=
\begin{pmatrix}
\text{analytic current-token score}
\\
\text{teacher suffix score at time }t
\end{pmatrix}
\\[1mm]
&\Downarrow
\\[-1mm]
&\score_h^\star(z,t\mid x)
=
\nabla_z\log m_{h,t}^{\targetpi}(z\mid x)
\\[1mm]
&\Downarrow
\\[-1mm]
&Z_T\sim\Normal(0,I)
\xrightarrow{\text{one reverse diffusion}}
W_H\sim\mathbb M_H^{\targetpi}
\\[1mm]
&\Downarrow
\\[-1mm]
&\E\left[\sum_{t<L_H(W_H)}R_t\right]
=
\E_{\targetpi}\left[\sum_{t=0}^{H-1}\gamma^tR_t\right].
\end{aligned}
}
\label{eq:complete-chain}
\end{equation}

\appendix

\section{Derivation from the Fokker-Planck Equation}
\label{app:fokker-planck}

For completeness, consider two diffusion processes with a shared diffusion
coefficient $g(t)$ and marginal densities $p_t^0,p_t^1$:
\begin{equation}
\dd Z_t=b_t^i(Z_t)\dd t+g(t)\dd W_t,
\qquad i\in\{0,1\}.
\end{equation}
Their Fokker-Planck equations are
\begin{equation}
\partial_t p_t^i
=-\nabla\cdot(p_t^i b_t^i)
+\frac{g(t)^2}{2}\Delta p_t^i.
\end{equation}
For $p_t=(1-\gamma)p_t^0+\gamma p_t^1$,
\begin{equation}
\partial_t p_t
=-\nabla\cdot
\left(
(1-\gamma)p_t^0b_t^0
+\gamma p_t^1b_t^1
\right)
+\frac{g(t)^2}{2}\Delta p_t.
\end{equation}
Hence the mixture path is generated by the drift
\begin{equation}
b_t(z)
=
\frac{
(1-\gamma)p_t^0(z)b_t^0(z)
+\gamma p_t^1(z)b_t^1(z)
}{
(1-\gamma)p_t^0(z)+\gamma p_t^1(z)
}.
\end{equation}
For reverse diffusions with the same forward drift and diffusion coefficient,
this reduces to the mixture-score formula in \cref{lem:mixture-score}. This is
the diffusion mechanism behind TD$^2$-DD and behind
\cref{eq:main-loss}.

\section{Detailed Proof of the Continuation Score}
\label{app:continuation-score}

Fix $x$ and write $x'=\nextx(y)$. Conditional on $Y=y$, the continuation
marginal is
\begin{equation}
p_t(z_0,z_+\mid y,x)
=
q_{t\mid0}(z_0\mid y)
 m_{h-1,t}^{\targetpi}(z_+\mid x').
\end{equation}
Therefore,
\begin{equation}
\nabla_{(z_0,z_+)}\log p_t(z_0,z_+\mid y,x)
=
\begin{pmatrix}
\nabla_{z_0}\log q_{t\mid0}(z_0\mid y)
\\
\nabla_{z_+}\log m_{h-1,t}^{\targetpi}(z_+\mid x')
\end{pmatrix}.
\end{equation}
Let
\begin{equation}
p_t^{\branchcont}(z\mid x)
=
\int\kappa_{\targetpi}(\dd y\mid x)
p_t(z\mid y,x).
\end{equation}
The conditional score identity gives
\begin{equation}
\nabla_z\log p_t^{\branchcont}(z\mid x)
=
\E
\left[
\nabla_z\log p_t(z\mid Y,x)
\middle|
Z=z,X=x,\text{continuation}
\right].
\end{equation}
Thus the target in \cref{eq:continuation-score-target} is an unbiased
conditional score target for the continuation marginal whenever the teacher is
exact.

\section{Discrete DDPM Parameterization}
\label{app:ddpm}

For a discrete DDPM schedule, let
\begin{equation}
Z_t
=
\sqrt{\bar\alpha_t}\,\Phi_h(W)
+
\sqrt{1-\bar\alpha_t}\,\epsilon,
\qquad
\epsilon\sim\Normal(0,I).
\end{equation}
The analytic token score is
\begin{equation}
\nabla_{z_t}
\log q_t(z_t\mid W)
=
-
\frac{
 z_t-\sqrt{\bar\alpha_t}\,\Phi_h(W)
}{1-\bar\alpha_t}.
\end{equation}
Equivalently, a noise-prediction network may be trained. On the stopping branch,
the target is the sampled Gaussian noise. On the continuation branch, the
current-token target is its sampled noise, while the suffix target is the
teacher's noise prediction at the same discrete time index. The resulting loss
is the trajectory counterpart of the TD$^2$-DD implementation in
\citet{farebrother2025tdflow}.

\subsection{Velocity Parameterization at the Exact Source}
\label{app:vpred}

The boundary condition $\alpha_T=0,\sigma_T=1$ in \cref{eq:boundary-schedule}
makes the source marginal exactly $\Normal(0,I)$, which is required for the
common-source statement in \cref{thm:path-mixture}. Under this exact boundary
the noise parameterization is numerically singular: recovering the clean tensor
from a predicted noise,
\begin{equation}
\widehat{\Phi_h(W)}
=
\frac{z_t-\sqrt{1-\bar\alpha_t}\,\hat\epsilon}{\sqrt{\bar\alpha_t}},
\end{equation}
divides by $\sqrt{\bar\alpha_T}=0$ at the terminal step, so any error in
$\hat\epsilon$ is amplified without bound. We therefore parameterize the network
by the velocity
\begin{equation}
v_t
:=
\sqrt{\bar\alpha_t}\,\epsilon
-
\sqrt{1-\bar\alpha_t}\,\Phi_h(W),
\end{equation}
whose clean-tensor recovery,
\begin{equation}
\widehat{\Phi_h(W)}
=
\sqrt{\bar\alpha_t}\,z_t
-
\sqrt{1-\bar\alpha_t}\,v_t,
\end{equation}
stays finite at the source, where it reduces to
$\widehat{\Phi_h(W)}=-v_T$. The velocity, the noise, and the score are related
by the invertible affine maps
\begin{equation}
\epsilon
=
\sqrt{\bar\alpha_t}\,v_t+\sqrt{1-\bar\alpha_t}\,z_t,
\qquad
\nabla_{z_t}\log q_t(z_t\mid W)
=
-\frac{\epsilon}{\sqrt{1-\bar\alpha_t}},
\end{equation}
so the velocity network encodes exactly the same score field
$\nabla_z\log m_{h,t}^{\targetpi}$ as
\cref{thm:population-score-backup}; the Bellman decomposition
\cref{eq:noise-bellman} and all recovery guarantees are unchanged. The branch
targets are the images of the $\epsilon$-targets under this map: on the stopping
branch the target is $v_t$ of the padded stop path
$\stoptraj_h(y)$; on the continuation branch the current-token target is the
analytic $v_t$ of the known transition $y$, and the suffix target is the frozen
teacher's velocity prediction at the same discrete index, again with the
stop-gradient of \cref{eq:continuation-score-target} and without any
$\gamma$ prefactor. This is the parameterization used by the reference
implementation, following the zero-terminal-SNR analysis of
\citet{lin2024commondiffusion} and the velocity objective of
\citet{salimans2022progressive}.

\section{Practical Variant with EMA Targets}
\label{app:ema}

The exact analysis freezes a horizon-$(h-1)$ teacher while fitting horizon $h$.
A practical implementation may use one shared network and an EMA target. Let
$\theta$ be the online parameters and $\bar\theta$ the target parameters. A
curriculum distribution $p(h)$ first concentrates on short horizons and then
expands toward $H$. For a sampled horizon $h$, the continuation target queries
$\score_{\bar\theta}(z_+,t\mid x',h-1)$. After each update,
\begin{equation}
\bar\theta
\leftarrow
\zeta\bar\theta+(1-\zeta)\theta.
\end{equation}
This replaces exact stage boundaries by a two-timescale approximation. The
stagewise theorem remains the clean population reference point.

\begin{thebibliography}{9}

\bibitem[Farebrother et al.(2025)]{farebrother2025tdflow}
Jesse Farebrother, Matteo Pirotta, Andrea Tirinzoni, R\'emi Munos,
Alessandro Lazaric, and Ahmed Touati.
\newblock Temporal Difference Flows.
\newblock In \emph{Proceedings of the 42nd International Conference on
Machine Learning}, 2025.

\bibitem[Ho et al.(2020)]{ho2020ddpm}
Jonathan Ho, Ajay Jain, and Pieter Abbeel.
\newblock Denoising Diffusion Probabilistic Models.
\newblock In \emph{Advances in Neural Information Processing Systems}, 2020.

\bibitem[Janner et al.(2020)]{janner2020gamma}
Michael Janner, Igor Mordatch, and Sergey Levine.
\newblock Gamma-Models: Generative Temporal Difference Learning for
Infinite-Horizon Prediction.
\newblock In \emph{Advances in Neural Information Processing Systems}, 2020.

\bibitem[Lin et al.(2024)]{lin2024commondiffusion}
Shanchuan Lin, Bingchen Liu, Jiashi Li, and Xiao Yang.
\newblock Common Diffusion Noise Schedules and Sample Steps are Flawed.
\newblock In \emph{IEEE/CVF Winter Conference on Applications of Computer
Vision (WACV)}, 2024.

\bibitem[Salimans and Ho(2022)]{salimans2022progressive}
Tim Salimans and Jonathan Ho.
\newblock Progressive Distillation for Fast Sampling of Diffusion Models.
\newblock In \emph{International Conference on Learning Representations}, 2022.

\bibitem[Song et al.(2021)]{song2021sde}
Yang Song, Jascha Sohl-Dickstein, Diederik P. Kingma, Abhishek Kumar,
Stefano Ermon, and Ben Poole.
\newblock Score-Based Generative Modeling through Stochastic Differential
Equations.
\newblock In \emph{International Conference on Learning Representations},
2021.

\bibitem[Vincent(2011)]{vincent2011score}
Pascal Vincent.
\newblock A Connection Between Score Matching and Denoising Autoencoders.
\newblock \emph{Neural Computation}, 23(7):1661--1674, 2011.

\bibitem[Ying et al.(2026)]{ying2026tddiffusion}
Qizhen Ying, Yangchen Pan, Victor Adrian Prisacariu, and Junfeng Wen.
\newblock Temporal Difference Learning for Diffusion Models.
\newblock In \emph{Proceedings of the 43rd International Conference on
Machine Learning}, 2026.

\end{thebibliography}

\end{document}
